\chapter{Registers and integer representation}
\MARSChapterHeading{Registers and integer representation}
\label{revision:integer-representation}

In the first program, we treated a register as a place to keep a number. That
model was sufficient for small calculations. Now consider a debugger showing
\texttt{0x0000000000000080}. Is the value 128 or $-128$? To answer, we need to
know which bits belong to the value and how those bits are interpreted.

We will write a small byte inspector. It accepts a decimal integer from 0 through
255 and prints two interpretations of its low eight bits: an unsigned value and
a signed value. Input 128 will produce 128 and $-128$. The program will help us
separate three questions: what bits are present, how wide the value is, and what
number those bits represent. Those questions also arise when a program reads a
pixel, a character or a sample from memory.

\section{Counting with bits}

A \emph{bit} has two possible values, zero and one. A group of eight bits is a
\emph{byte}. Binary notation uses powers of two as place values, just as decimal
notation uses powers of ten. In the byte \texttt{10010110}, the set bits contribute
$128+16+4+2=150$ when we read the pattern as an unsigned number.

\begin{center}
\begin{tabular}{lrrrrrrrr}
\toprule
Bit position & 7 & 6 & 5 & 4 & 3 & 2 & 1 & 0 \\
Place value & 128 & 64 & 32 & 16 & 8 & 4 & 2 & 1 \\
Example bit & 1 & 0 & 0 & 1 & 0 & 1 & 1 & 0 \\
\bottomrule
\end{tabular}
\end{center}

Bit positions count from zero at the right. We call bit 0 the \emph{least
significant bit} because changing it changes an unsigned value by only one.
Bit 7 is the most significant bit of this byte. Its place value is 128.

Eight bits have $2^8=256$ distinct patterns. An unsigned interpretation assigns
them the integers 0 through 255. The largest pattern, \texttt{11111111}, represents
$128+64+32+16+8+4+2+1=255$. There is no ninth position available to store 256.
Width therefore limits the values a representation can hold.

A MIPS64 general-purpose integer register has 64 bits. We can use its low eight
bits to hold a byte value, but doing so does not shrink the register. The remaining
56 bits still exist. An instruction's rules determine whether those other bits
are kept, cleared or filled from a sign bit.

\section{Reading hexadecimal}

Writing 64 separate binary digits is cumbersome. \emph{Hexadecimal} notation groups
bits four at a time. Each group has sixteen possible patterns, written with the
digits 0 through 9 and the letters A through F. A represents ten; F represents
fifteen. In source code, the prefix \texttt{0x} identifies a hexadecimal constant.

For our example, split \texttt{10010110} into \texttt{1001} and \texttt{0110}.
Those groups represent 9 and 6, so the byte is \texttt{0x96}. Its decimal value is
$9\times16+6=150$. The notation changes; the underlying bits do not.
\texttt{0x80}, \texttt{10000000} in binary, and 128 in decimal all describe the
same unsigned byte value.

Two hexadecimal digits describe eight bits. Sixteen hexadecimal digits describe
64 bits. Leading zeros let the debugger show the full register width:
\texttt{0x0000000000000096} still has the unsigned value 150. When you inspect a
register, count the digits before deciding whether you are looking at a byte,
a 32-bit word or a complete 64-bit value.

\section{Giving the same bits a signed meaning}

Programs also need negative numbers. The signed integer representation used here
is \emph{two's complement}. For an eight-bit value, the top bit has weight $-128$
instead of $+128$; the other seven weights are unchanged. Thus
\texttt{10010110} represents $-128+16+4+2=-106$ as a signed byte.

An equivalent calculation is useful when the unsigned value is already known.
If bit 7 is zero, the signed and unsigned interpretations agree. If bit 7 is one,
subtract 256 from the unsigned value. Our unsigned 150 becomes $150-256=-106$.
This gives the signed range $-128$ through 127.

\begin{center}
\begin{tabular}{llll}
\toprule
Eight bits & Hexadecimal & Unsigned & Signed \\
\midrule
\texttt{00000000} & \texttt{0x00} & 0 & 0 \\
\texttt{01111111} & \texttt{0x7F} & 127 & 127 \\
\texttt{10000000} & \texttt{0x80} & 128 & $-128$ \\
\texttt{10010110} & \texttt{0x96} & 150 & $-106$ \\
\texttt{11111111} & \texttt{0xFF} & 255 & $-1$ \\
\bottomrule
\end{tabular}
\end{center}

The top bit is not a separate minus sign attached to a positive magnitude.
For example, \texttt{10000001} is $-128+1=-127$, rather than $-1$.
The all-one pattern represents $-1$ because $-128+127=-1$.

A register does not carry a label saying that its contents are a signed byte.
The instruction, service or program convention supplies that interpretation.
Before calling a debugger display wrong, ask which width and interpretation it
is using. A signed 64-bit display of \texttt{0x0000000000000080} correctly shows
128: bit 63 is zero. Reading just the low byte as signed gives a different answer.

\section{Extending a value to register width}

Suppose a byte's bits must become a value in a 64-bit register. For an unsigned
byte, we preserve its value by filling the upper bits with zeros. This operation
is called \emph{zero extension}. The byte \texttt{0x80} becomes
\texttt{0x0000000000000080}, still 128.

For a signed byte, preserving its value requires \emph{sign extension}: copy
bit 7 into every new upper position. A signed \texttt{0x80} becomes
\texttt{0xFFFFFFFFFFFFFF80}, whose signed 64-bit value is $-128$.
A signed \texttt{0x7F} becomes \texttt{0x000000000000007F}, or 127, because its
sign bit was zero. Sign extension preserves a signed value when widening it.
It does not preserve the unsigned interpretation of a pattern whose top bit is one.

The instruction \texttt{seb \$t1, \$t0} performs this byte sign extension. Its
name abbreviates ``sign-extend byte.'' It reads the low eight bits of
\texttt{\$t0} and writes the sign-extended result to \texttt{\$t1}. It neither
changes \texttt{\$t0} nor fetches a byte from memory. Here, ``byte'' specifies the
width of the source portion within a register.

We will obtain the unsigned version with \texttt{andi \$t0, \$v0, 0xff}.
The instruction performs a bitwise AND: each result bit is one only when both
corresponding input bits are one. The constant \texttt{0xff} has ones in its low
eight positions and zeros elsewhere after the immediate is zero-extended.
Consequently, the instruction keeps the source's low eight bits and clears all
higher bits. Such a constant is called a \emph{mask}. It selects the positions we
want to retain.

For input 150, the mask produces \texttt{0x0000000000000096}. The following
\texttt{seb} produces \texttt{0xFFFFFFFFFFFFFF96} in a different register.
We now have two full-width values that preserve the two meanings of one byte.
The masked source also satisfies \texttt{seb}'s requirement for a sign-extended
32-bit source value: its bit 31 and all higher bits are zero.

\section{A complete byte inspector}

The console interface is the one used in Chapter~\ref{revision:first-program}:
service 5 reads a decimal integer into \texttt{\$v0}, service 4 prints a string,
service 1 prints a signed 32-bit integer from \texttt{\$a0}, and service 10 exits.
Our two results lie within that printing service's range. A 64-bit register does
not make this particular service a general 64-bit integer printer.

\MARSCodeLabel{One byte, two integer interpretations}
\begin{minted}{mars-idea}
.data
prompt:   .asciiz "Enter an integer from 0 to 255: "
unsigned: .asciiz "Unsigned byte: "
signed:   .asciiz "Signed byte: "
newline:  .asciiz "\n"

.text
.globl main
main:
    la $a0, prompt
    li $v0, 4
    syscall
    li $v0, 5
    syscall
interpret:
    andi $t0, $v0, 0xff
    seb $t1, $t0
report:
    la $a0, unsigned
    li $v0, 4
    syscall
    move $a0, $t0
    li $v0, 1
    syscall
    la $a0, newline
    li $v0, 4
    syscall
    la $a0, signed
    li $v0, 4
    syscall
    move $a0, $t1
    li $v0, 1
    syscall
    la $a0, newline
    li $v0, 4
    syscall
    li $v0, 10
    syscall
\end{minted}

Save the source as \texttt{byte-inspector.asm}. Select the MIPS64 Release 6 target,
run it, and enter 150. The console should print \texttt{Unsigned byte: 150} and
\texttt{Signed byte: -106} on separate lines. The prompt and input may share a
line in the IDE console. The values are held in \texttt{\$t0} and \texttt{\$t1}
while \texttt{\$v0} is reused for service selectors.

The input range is a precondition for this exercise, not a check performed by
the program. Entering 256 will leave zero after the mask; it will not produce an
error message. Masking selects bits. It does not validate that an input belongs
to a stated range. We will add explicit input decisions after learning branches.

Before another run, predict the two output numbers for 0, 127, 128 and 255.
Use the preceding table to check your reasoning, then compare your predictions
with the console. These cases distinguish the ends of the range from the point
where the signed interpretation changes.

\section{Inspecting the conversion}

Set a breakpoint at \texttt{interpret} and start Debug with input 128. Before
\texttt{andi} executes, \texttt{\$v0} contains the input. Step once: the complete
\texttt{\$t0} should be \texttt{0x0000000000000080}. Step again: the complete
\texttt{\$t1} should be \texttt{0xFFFFFFFFFFFFFF80}. In a signed 64-bit decimal
view those registers show 128 and $-128$. Inspect the full values, since their
last two hexadecimal digits are identical.

If both output lines say 128, check two boundaries. Did \texttt{seb} write
\texttt{\$t1}, and did the second integer request copy \texttt{\$t1} into
\texttt{\$a0}? A correct conversion followed by a copy from \texttt{\$t0} still
prints the unsigned result. Changing the debugger's display format cannot repair
a wrong argument in the running program.

For another experiment, replace \texttt{seb \$t1, \$t0} with
\texttt{move \$t1, \$t0}. Predict which inputs will expose the change.
Inputs below 128 cannot distinguish these programs because their sign bit is zero.
Input 128 can. A useful test exercises the property on which two implementations
differ; a successful run with zero alone provides little evidence about sign extension.

\section{Why instruction width matters}

The first chapter used \texttt{addu} and \texttt{addiu}. These instructions compute
32-bit word results, then sign-extend those results into the 64-bit destination.
The \texttt{u} means that signed overflow does not raise an arithmetic exception.
It does not request zero extension.
Their sources must themselves be sign-extended 32-bit values. The small inputs
in Chapter~\ref{revision:first-program} and the source below satisfy that rule.

Consider the following short experiment. Stop before the exit service and inspect
both result registers in hexadecimal and signed 64-bit decimal.

\par\noindent\begin{minipage}{\linewidth}
\MARSCodeLabel{Crossing the signed word boundary}
\begin{minted}{mars-idea}
.text
.globl main
main:
    li $t0, 0x7fffffff
    addiu $t1, $t0, 1
    daddiu $t2, $t0, 1
inspect:
    li $v0, 10
    syscall
\end{minted}
\end{minipage}\par

The initial value is $2^{31}-1=2147483647$. Adding one mathematically gives
2147483648. The word result of \texttt{addiu} has bit pattern
\texttt{0x80000000}. Sign extension writes \texttt{0xFFFFFFFF80000000} into
\texttt{\$t1}; its signed 64-bit interpretation is $-2147483648$.
\texttt{daddiu} performs a 64-bit addition instead. It writes
\texttt{0x0000000080000000} into \texttt{\$t2}, whose signed value is 2147483648.
Both instructions receive the same source and immediate. Their arithmetic widths
explain the different destinations.

The immediate in both instructions is a signed 16-bit constant extended to the
operation's width. The \texttt{u} does not make that immediate unsigned either.
The constant one is representable, so this detail does not change our example.
When selecting an instruction, identify its operand rules and result width;
register width alone is insufficient.

\section{Writing and explaining a variation}

Extend the byte inspector to print the unsigned value plus one on a third line.
Keep that sum as an ordinary register value: for input 255 it should print 256.
Choose where the sum must survive while strings and other numbers are printed.
Test 0, 127, 128 and 255, and explain which instruction supplies each output argument.

Then change the third result to retain only the low eight bits of that sum.
For input 255 the new result should be zero. This is \emph{wrapping} to eight bits:
results repeat every 256 values. Identify the instruction that discards the extra
bit. Distinguish this deliberate narrowing from the sign extension that made the
signed output negative.

Finally, explain why printing a value with service 1 is an insufficient test of
all 64 register bits. That service consumes a signed 32-bit integer. The debugger's
full-width view is needed for the preceding word-boundary experiment. The next
chapter will move byte values between registers and memory, where selecting the
right load width and extension becomes part of every access.
